\section{Native dual normalization with retained interfaces}
\label{sec:geometry}

\subsection{Source contract and size parameters}
The native normal-surface validator supplies a connected finite triangulation
of a compact orientable three-manifold with one torus boundary.  Generalized
face pairings and self-identifications are permitted only when the validator's
manifold checks pass.  The normal vector is nonnegative and integral, satisfies
the matching equations and quadrilateral constraints, and belongs to the
existing source contract.  We do not weaken any of these checks.

Write $t$ for the number of tetrahedra, and $B$ for the largest bit length of
a coordinate or local edge weight, with $B\geq1$.  The source constructor
\code{normal_parallelity_geometry.py} partitions the cut into interval-bundle
regions and a finite complementary handle core.  Its $i$-handles are denoted
$H_i$.  The core has boundedly many local nonparallel regions per source
tetrahedron.  Crude bounds adequate here are
\begin{equation}\label{eq:handle-bounds}
 |H_0|\leq6t,\quad |H_1|\leq8t,\quad |H_2|\leq12t,
 \quad |H_3|\leq4t.
\end{equation}
Indeed, a tetrahedron contributes at most two central pieces and four
triangle tips; a paired face contributes at most four nonparallel regions;
each interior source edge contributes at most its two extreme gaps; and
three-handles correspond to interior source vertices.  The total number of
one-handle crossings in all source two-handle cycles is at most $12t$:
each local tetrahedral edge occurs in at most two extreme cycles.

The encoded core is thus small even when a normal coordinate is enormous.
Its three-handles retain their source vertex and incident link triangles.
Each two-handle retains its source edge, its gap rank, and its full cyclic
attaching chart.  These records are geometric incidence data; they are
strictly stronger than a signed cellular boundary matrix.

\subsection{The graph of actual belt-point occurrences}
\begin{definition}[Rooted dual graph]
Let $K$ be an authenticated finite three-dimensional handle core.  The
\emph{dual graph} $D$ has one vertex for each three-handle and a formal root
$\rootv$ for boundary not capped by a three-handle.  A two-handle is
$D^2\times I$ and has belt sphere $\{0\}\times\partial I$, consisting of
two points.  Its edge joins the labels of these two actual belt-point
occurrences.  A point in a three-handle's attaching sphere receives that
vertex label; an uncapped point receives $\rootv$.
\end{definition}

The root need not be a connected boundary component.  It is an algebraic
device joining all places at which cancellation can terminate.  Edges from
the root to itself and edges whose two ends are the same three-handle are
loops.  A loop has two occurrences, not zero.  This distinction prevents a
signed boundary coefficient from being mistaken for geometric cancellation.

For native cores, the graph is recovered without tracing an expanded normal
surface.  Fix an oriented source edge $e$, with original endpoints $u,v$ and
normal intersection weight $w_e$.
\begin{enumerate}
\item If $w_e=0$, there is a single core gap.  Its dual edge joins the labels
  of $u$ and $v$.
\item If $w_e>0$, there are two extreme core gaps.  Rank zero joins $u$ to
  $\rootv$, and rank $w_e$ joins $\rootv$ to $v$.
\end{enumerate}
An original boundary vertex is assigned $\rootv$.  The other $w_e-1$ gaps
belong to the parallelity structure and are never enumerated by this step.

\begin{lemma}[Correct source orientation]\label{lem:orientation}
The native source edge identifier is a representative of the signed
union-find class of local tetrahedral edges.  At that representative the
orientation parity is zero.  Thus decoding its local edge $(a,b)$ with
$a<b$ identifies the original endpoint of rank zero as $a$, and that of the
maximum rank as $b$, even when the two global endpoint vertices coincide.
\end{lemma}
\begin{proof}
The parity stored by the signed equivalence structure is the parity between
a local edge and its representative.  A representative is related to itself
with parity zero.  The source rank convention reverses ranks precisely when
this parity is one.  At the representative no reversal occurs, so the
claimed identification follows directly from the local edge parameter.
The global equality of vertex labels does not identify the two local endpoint
germs.
\end{proof}

One must not use an unordered endpoint cache to choose between the two
extreme ranks.  The implementation decodes the representative itself; the
independent checker reconstructs the same geometric fact separately.

\begin{lemma}[Root reachability]\label{lem:rooted}
For every native source and admissible normal vector, every vertex of $D$
is connected to $\rootv$.
\end{lemma}
\begin{proof}
The one-skeleton of the connected source triangulation is connected and
contains boundary vertices.  Identify all boundary vertices with one root.
Boundary edges then become root loops; removing them does not affect root
reachability.  Keep every interior source edge of zero normal weight.  Replace
each positive-weight interior edge by two root stubs, one from each of its
original endpoints to the root.  Replacing an edge in this way preserves root
reachability: any path that formerly used that edge can pass through the root
instead.  The resulting graph, after discarding irrelevant root loops, is
exactly the graph just compiled from extreme core gaps.  Its nonroot vertices
are exactly the interior source vertices, equivalently the core three-handles.
\end{proof}

\subsection{Geometric cancellation, including the relative assertion}
The following is a source-specific use of the classical handle cancellation
theorem: the attaching sphere of a $(k+1)$-handle must meet the belt sphere
of a $k$-handle transversely in one point \cite[Theorem~5.4]{Milnor1965}.
Here the belt sphere is zero-dimensional, so the compiled occurrence data
give the required geometric intersection directly.

\begin{lemma}[Rooted dual forest cancellation]\label{lem:dual-cancel}
Suppose every component of an authenticated dual graph meets its root.
For any spanning tree $T$ after the roots have been identified, there is a
sequence of geometric two-/three-handle cancellations deleting precisely
the two-handles labeled by $T$ and every three-handle.  The surviving
two-handles retain their attaching maps to the abstract zero-/one-handle
base.  No algebraic free reduction or two-handle slide is needed.

If a specified exposed boundary subcomplex has a collar disjoint from the
dual cancellation balls and tubes, the reduction preserves its labeled
embedding up to the relative boundary identification.
\end{lemma}
\begin{proof}
Orient the tree toward the root.  Choose a vertex adjacent to the already
exposed root region.  Its parent edge has one belt point in that
three-handle's attaching sphere and the other in exposed boundary.  The
classical cancellation criterion therefore applies.  Cancel the pair.
Other edges formerly incident with that three-handle now meet exposed
boundary.  On the dual graph this is precisely contraction of the selected
edge into the root.  Continue outward along the tree.  Each step removes one
three-handle; after $|H_3|$ steps the selected tree edges and all three-handles
have disappeared.

There is an elementary boundary-ball proof of each step, which also
establishes the relative assertion.  Write the cancelling two-handle as
$h_e=D^2\times[-1,1]$, with caps $D_e^\pm$.  Label it so that the
three-handle's attaching sphere $\Sigma$ contains $D_e^+$ and not
$D_e^-$.  Let $R$ contain all other handles.  Other three-handle spheres
miss both caps, so their attachments can be made before this pair.  Then
$\Delta=\overline{\Sigma\setminus D_e^+}$ is a disk in $\partial R$.
Together with the attaching annulus $A_e$ of $h_e$, it forms a disk
$E=\Delta\cup A_e$ in $\partial R$.  Meanwhile
$h_e\cup_{D_e^+}h_3$ is a three-ball $B$, and the original manifold is
$R\cup_E B$.  Adding a ball along a boundary disk adds only a collar.
The collar can be straightened relative to its boundary circle and to
every exposed patch in $R$ outside the interior of $E$.

The move alters the upper attachment boundary but does not introduce a
slide of any retained two-handle over another one.  Consequently the
attaching maps of all retained two-handles to the abstract lower base can
be kept as before.  This is a statement about the resulting handle
description.  It does not assert that a homeomorphism fixes every point of
the entire old lower submanifold.
\end{proof}

\begin{lemma}[Native exposed ports]\label{lem:ports}
The exposed core-facing ports of the native parallelity construction can be
retained by the reduction in Lemma~\ref{lem:dual-cancel}, with their incidence,
cyclic order, and side labels unchanged.
\end{lemma}
\begin{proof}
The native constructor places each core-facing port on an exposed region of
a zero-handle or a one-handle; its chart records \code{core_dimension} as
zero or one.  There are no ports in the interiors of core two- or
three-handles.  They are actual boundary patches of the complementary core,
not patches later covered by a three-handle.

In the boundary-ball proof of Lemma~\ref{lem:dual-cancel}, every port is
therefore contained in $R$.  It misses the interior of $E$: that interior
is an internal gluing region in $R\cup_E B$, not exposed boundary.
If a port edge meets $\partial E$, use the collar straightening fixed on
that circle.  The relative identification thus fixes all exposed port
patches.  Iterating this argument retains their product-chart boundary
identifications.  No port is deleted, reordered, or reattached, so its
signed word and its incidence with the unchanged abstract lower charts
are retained.
\end{proof}

This argument applies to the authenticated native ports.  A proposed future
boundary pattern that crosses the cancellation support must supply its own
transport.  The mere fact that a pattern has a small graph or an Euler
characteristic does not meet this requirement.

\begin{theorem}[Native top-handle normalization]\label{thm:native}
Given an authenticated native parallelity geometry, all core three-handles
can be removed by a certified geometric reduction.  Its reduced core and
retained exposed interfaces describe the same cut manifold with the same
parallelity attachments.  The reduction removes exactly $s=|H_3|$
two-handles, preserves Euler characteristic, and needs $O(t)$ graph records.

Graph compilation and a breadth-first tree use $O(t)$ arithmetic/label
operations.  Independent union-find verification uses
$O(t\alpha(t))$ label operations.  The arithmetic involves $B$-bit weights
and $O(\log(t+1))$-bit labels, so a deterministic array implementation has
polynomial bit cost, for example
\[
 O\bigl(tB+t\alpha(t)\log(t+1)\bigr)
\]
in a logarithmic-cost random-access model, after source authentication,
for the canonical bounded-size certificate.  The delivered producer uses
arrays and breadth-first search; replay uses separate arrays and union-find.
A multi-tape Turing implementation also has polynomial bit cost, without
asserting this exact random-access bound.  For arbitrary malformed input,
verification time additionally includes its actual encoded length; a huge
invalid identifier must be read before it can be rejected.
\end{theorem}
\begin{proof}
Apply Lemma~\ref{lem:rooted}, construct a tree, and use
Lemmas~\ref{lem:dual-cancel} and~\ref{lem:ports}.  A tree on $s+1$ vertices
has $s$ edges.  Removing equal numbers of two- and three-handles leaves the
alternating handle count unchanged.  The size bounds
\eqref{eq:handle-bounds} and the local edge-weight formula give the operation
bounds.  The independent checker needs only to reconstruct the selected
edge endpoints, reject cycles, and verify that all $s$ vertices reach the
root.  The list of selected original two-handle identifiers has
$O(s\log(t+1))$ bits; all retained charts can be shared or copied from the
authenticated source.
\end{proof}

\begin{corollary}[Support-stable normalization]\label{cor:support}
For a fixed source triangulation, the dual graph, using semantic labels
$(e,\text{zero})$, $(e,\text{left})$, and $(e,\text{right})$, depends only
on the set of edges with zero normal weight.  A selected dual tree is valid
throughout that stratum.  In particular it is valid for every positive
integer multiple of a supplied normal vector.
\end{corollary}
\begin{proof}
The two cases in the graph compiler use only $w_e=0$ or $w_e>0$.
The endpoint labels come from the source vertices and remain fixed.
Multiplication by a positive integer preserves every zero test.  Maximum
ranks are updated to the current $w_e$ when actual chart identifiers are
reconstructed.
\end{proof}

\begin{proposition}[Optimal initial deletion cost]\label{prop:max-tree}
Assign a nonnegative cost $c_e$ to every native two-handle, such as its
attaching-cycle length or its raw occurrence count after a fixed lower
forest.  A maximum-weight spanning tree minimizes the total cost of the
retained two-handles among all dual-tree normalizations.
\end{proposition}
\begin{proof}
For a tree $T$, the retained cost is
$\sum_{e\in E(D)}c_e-\sum_{e\in T}c_e$.  The first summand is constant.
Thus minimizing the retained cost is exactly the maximum spanning-tree
problem.
\end{proof}
This objective is inexpensive but does not guarantee subsequent raw
cancellability.  The next sections address that stricter question.
